\section{원본 \bench 검증}
\label{app:verify_tau}

\subsection{구현 검증}

구현 오류를 최소화하기 위해, 우리는 (1) 모든 도메인에 걸쳐 통합된 도구 형식을 제공하여 에이전트 역량이 표현되는 방식의 일관성을 보장하고, (2) 각 도메인 환경이 데이터 모델을 명시하도록 보장하며, (3) 단위 테스트를 위해 특별히 설계된 모의 도메인을 도입하여 핵심 벤치마크 기능을 격리해 검증할 수 있게 했다.

\subsection{과제 검증}

과제의 명확성과 정확성은 의미 있는 평가에 필수적이다. 우리는 여러 조치를 통해 과제 명세를 개선했다:
\begin{itemize}
    \item \textbf{구조화된 과제 데이터:} 각 과제의 \textbf{목적}(즉, 테스트되는 구체적 역량)을 자세히 설명하는 메타데이터로 과제 설명을 보강했다. 사용자 지시는 의도, 구체적 지시사항, 알려진/알 수 없는 정보와 같은 차원을 따라 구조화했다. 공통 데이터베이스 기본 상태 위에 각 과제를 구축할 수 있는 선택지를 유지하면서도 시작 상태를 더 세밀하게 제어할 수 있도록 \textbf{초기화 옵션}을 도입했다.
    \item \textbf{분류된 평가:} 과제 평가 기준을 보강하고 서로 다른 범주로 나누었다: \textbf{환경 단언}(예: 데이터베이스 상태 확인), \textbf{커뮤니케이션 단언}(에이전트가 전달한 정보 검증), \textbf{자연어 단언}(디버깅을 쉽게 하기 위해 자연어로 지정된 세밀한 확인 허용), \textbf{행동 단언}(필수 에이전트 행동 확인).
    \item \textbf{반복적 검토 프로세스:} 시뮬레이션 결과에 기반한 반복적 검토 프로세스를 구현했다. 각 과제에 대해 시뮬레이션을 실행한다. 검토자는 일시적인 에이전트 또는 사용자 시뮬레이터 오류가 시뮬레이션을 조기에 중단시킬 수 있는 경우 개입해 이를 수정할 수 있으며, 이로써 과제를 완전히 탐색할 수 있다. 그런 다음 시뮬레이션 결과를 검토하여 \textbf{과소명세}, \textbf{과잉명세}, \textbf{비고유 해법} 같은 문제가 있는지 확인한다. 검토 결과에 따라 과제 지시사항을 개선한다.
    \item \textbf{프로그램 방식 과제 생성:} 새로 도입한 도메인에서는 \textbf{프로그램 방식 과제 생성}을 \textbf{자동 검증}과 결합해 사용하여 설계상 정확성을 보장한다.
\end{itemize}
